\chapter{Binary search trees}\label{chap:binary:search:trees}

Binary search trees store lower (\textit{resp.} higher) keys than the root in the left (\textit{resp.} right) subtree. Searches therefore are expected to occur with logarithmic time complexity and implementation is pretty straightforward. These structures natively allow duplicate entries so be careful when using them for implementing maps.

The main issue with binary search trees is to keep them as balanced as possible so that searches keep as logarithmic as possible (even in the case of worst-case, degenerate entries). One (complicated) method known as Red-Black trees will ensure balanced trees at all times with hard worst-case time guarantees.

One slightly less complicated and even more beautiful structure is called Splay trees. In addition to removals, Splay trees also get modified during lookups, thereby implementing a form of Move-to-Front optimization like in Chap.~\ref{chap:hashtables:closed}. Since the root of a Splay tree actually \textit{moves} during searches, you definitely need a header for your tree data structures to enforce our coding rules---they are more tedious to get right than the treaps below.

A treap builds a heap structure on top of a binary search tree (hence the portemanteau in the name). This is called a Cartesian tree~\cite{ACP:treaps} and it uses hashing to implement randomized, balanced search trees (the hash provides the required source of randomness). It is obviously a beautiful data structure, yet certainly it cannot be the first you implement---more likely the second!

One other great variant that must be mentioned is called B-trees, which minimize the amount of data to be moved around during insertions/removals: this is particularly well suited for large, offline databases (on disk, like for implementing a filesystem!) B-trees are arguably the trickiest of the four to get to work and, given their focus on out-of-core databases, the need for generic B-trees is quite unlikely.

Remove these paragraphs when you're done!

\section{Design rationale}

\section{Binary search tree interface}

\section{Example variant: Treaps}

\section{Example variant: Splay trees}

